%----------------------------------------------------------------------
\section{Worked example: $\kappa$-equivalence batched across a population}
\label{app:feedback-loop}
%----------------------------------------------------------------------

Appendix \ref{app:examples} grounded the planner's
arithmetic on a simple workload with MODEL-4's stated
parameters. This appendix extends the exercise to a
non-trivial, parameterized algorithm from the existing
CSTZ codebase and traces the design-space dynamics the
planner exposes. The point is no longer ``see, the
planner makes this fast''; it is ``see, the planner is
a feedback loop in the engineering search space.''

\subsection{The algorithm}\label{app-f:alg}

$\kappa$-equivalence on pairs of elements,
\texttt{cstz.sets.kappa\_equiv} (\S[DAF] 4.3, Def
\ref{def:exec-operationalist}), is the fundamental
operationalist check. Batched across a population of
$K$ elements under a regime of $M$ discriminators at
dimension $n$, the algorithm computes the $K \times K$
equivalence matrix $E \in \GF(2)^{K \times K}$:
\begin{equation*}
  E[i, j] = \bigwedge_{m=1}^{M} \bigl[\langle d_m, a_i \rangle = \langle d_m, a_j \rangle\bigr].
\end{equation*}
In [DAF]'s language this is the sheafification of the
local admissibility-match predicate over all ordered
pairs; operationally it is the core of the
\texttt{observe.ObservationState} equivalence-class
computation.

The free parameters are $(K, n, M)$. Realistic values:
$K \in \{2^3, 2^8, 2^{12}\}$, $n \in \{3, 8, 16\}$,
$M \in \{3, n, 2n\}$. The planner's job is to choose
an execution strategy that fits a stated hardware
profile; we will see that the ``right'' strategy is
not a function of $(K, n, M)$ alone but of $(K, n, M)$
against a budget vector $B$.

\subsection{Three plan families}\label{app-f:families}

Each family can be further refined (fused vs.\
time-shared, $\Lambda$ vs.\ $\mathrm{RM}$ basis) but
the high-level shapes are:

\paragraph{Plan $\alpha$: pairwise na\"{\i}ve.}
Iterate over $\binom{K}{2}$ pairs; for each, run
\texttt{kappa\_equiv} synchronously on a single lane.
$\VRAM$ traffic is
$\binom{K}{2} \cdot (2n + \lceil M/8 \rceil)$ bytes.
$\LAT$ is $\binom{K}{2}$ iterations at the
single-lane cost of one \texttt{kappa\_equiv}
evaluation ($M$ AND-reductions of $n$-bit dot
products). No batching advantage; minimum $\REG$
(two elements live).

\paragraph{Plan $\beta$: dot-product batching in $\Lambda$.}
For each discriminator $d_m$, batch-compute $(d_m \cdot a_i)_{i=1}^K$
in a single warp pass of $\lceil K/w \rceil$
iterations. This populates a
$K \times M$ \emph{evaluation matrix} $V$ with
$V[i, m] = \langle d_m, a_i \rangle$. The equivalence
matrix is then $E[i, j] = 1 - \max_m |V[i, m] -
V[j, m]|$, computed by XOR-reducing $M$-bit rows
pairwise. $\VRAM$ writes: $KM$ bits of $V$.
$\LAT$: $M \cdot \lceil K/w \rceil$ dot-product
iterations plus $\binom{K}{2} \cdot \lceil M/w \rceil$
XOR-reduction iterations. $\REG$: moderate
($V$'s row width plus accumulators).

\paragraph{Plan $\gamma$: encode to $\mathrm{RM}$, syndrome-compare.}
Butterfly-encode each $a_i$ into $\zeta(a_i) \in
\GF(2)^{2^n}$ (Section \ref{sec:encoding}). Under the
RM basis, $\kappa$-equivalence at the discriminator
set $\kappa$ reduces to: $\zeta(a_i) \oplus \zeta(a_j)$
has zero syndrome under the parity-check matrix
obtained from $\kappa$. Encode cost: $K \cdot n \cdot
\lceil 2^n / w \rceil$ XORs. Pairwise comparison:
$\binom{K}{2} \cdot \lceil 2^n / w \rceil$ XORs plus
syndrome GEMV. $\REG$: high ($2^n$-byte registers).

\subsection{Concretization under MODEL-4}
\label{app-f:model4}

Using the MODEL-4 parameters of Appendix
\ref{app-e:platform} ($w = 8$, $\alpha = 4$, $\beta =
1$, $\gamma = 4$, $\ell_\oplus = 1$), we derive
exact cost vectors for each family at three
workload instances. All numbers are arithmetic
consequences of the parameters above; none are
measured.

\paragraph{Instance 1: small population.}
$(K, n, M) = (8, 3, 3)$. $\binom{K}{2} = 28$.
\begin{center}
\begin{tabular}{l|rrrr}
Plan & $\VRAM$ & $\SMEM$ & $\REG$ & $\LAT$ \\
\hline
$\alpha$ & $28 \cdot (6 + 1) = 196$ & 0 & 2 & $\gamma + 28 \cdot (M \cdot \ell_\oplus + 1) = 116$ \\
$\beta$ & $KM + K \cdot n = 48$ & 0 & 3 & $\gamma + M + 28 \cdot 1 = 35$ \\
$\gamma$ & $K \cdot 2^n + K \cdot 2^n / w = 72$ & 0 & 8 & $\gamma + n \cdot K + 28 \cdot (2^n + M) = 332$ \\
\end{tabular}
\end{center}
Frontier: $\{\beta\}$. Plan $\beta$ dominates both
alternatives on every axis except $\REG$ (where
$\alpha$ is smaller by 1); at any $B_{\REG} \geq
3$, $\beta$ is the sole frontier element. Plan
$\gamma$'s encoding overhead swamps its syndrome
advantage at this small $n$.

\paragraph{Instance 2: moderate population, larger $n$.}
$(K, n, M) = (256, 8, 8)$. $\binom{K}{2} = 32{,}640$.
\begin{center}
\begin{tabular}{l|rrrr}
Plan & $\VRAM$ & $\SMEM$ & $\REG$ & $\LAT$ \\
\hline
$\alpha$ & $32{,}640 \cdot 17 \approx 5.5 \cdot 10^5$ & 0 & 2 & $\gamma + 32{,}640 \cdot 9 \approx 2.9 \cdot 10^5$ \\
$\beta$ & $KM + Kn = 4096$ & 0 & 3 & $\gamma + 8 \cdot 32 + 32{,}640 \cdot 1 \approx 3.3 \cdot 10^4$ \\
$\gamma$ & $K \cdot 2^8 + \cdots \approx 73{,}728$ & 256 & 32 & $\gamma + 8 \cdot 256 + 32{,}640 \cdot 40 \approx 1.3 \cdot 10^6$ \\
\end{tabular}
\end{center}
Frontier still $\{\beta\}$ under generous budgets;
under a tight $\VRAM$ budget ($B_{\VRAM} < 4096$),
all three are infeasible and the workload cannot run.
$\gamma$'s growth in $2^n$ is prohibitive.

\paragraph{Instance 3: large $n$, small $M$.}
$(K, n, M) = (256, 16, 4)$. Here the
grade-structure of $\kappa$ (only 4 of 16 possible
discriminators active) means $\gamma$ can use a
\emph{grade-truncated} RM basis of dimension
$R_1 = 1 + n = 17$ rather than $2^{16} = 65{,}536$.
\begin{center}
\begin{tabular}{l|rrrr}
Plan & $\VRAM$ & $\SMEM$ & $\REG$ & $\LAT$ \\
\hline
$\alpha$ & $32{,}640 \cdot 33 \approx 1.1 \cdot 10^6$ & 0 & 2 & $\gamma + 32{,}640 \cdot 17 \approx 5.5 \cdot 10^5$ \\
$\beta$ & $KM + Kn = 5120$ & 0 & 3 & $\gamma + 4 \cdot 512 + 32{,}640 \approx 3.4 \cdot 10^4$ \\
$\gamma^*$ (grade $\leq 1$) & $K R_1 + K R_1 / w = 4896$ & 0 & 5 & $\gamma + n \cdot K + 32{,}640 \cdot (R_1 + M) \approx 6.9 \cdot 10^5$ \\
\end{tabular}
\end{center}
Frontier: $\{\beta\}$ for $\LAT$-binding budgets;
$\gamma^*$ is a close second and becomes frontier when
$B_{\REG} \leq 3$ is relaxed to allow $R_1$-bit
registers. The grade truncation is the key lever:
without it, $\gamma$ is infeasible at $n = 16$; with it,
$\gamma^*$ is competitive.

\subsection{How complexity questions are concretized}
\label{app-f:complexity}

Three classical complexity questions, restated in the
planner's vocabulary:

\emph{(a) Is this algorithm dominated by compute or by
memory?} Plan $\beta$'s $\LAT$ in Instance 2 is
$3.3 \cdot 10^4$ cycles; if that exceeds $B_{\LAT}$,
the binding resource is $\LAT$ and the dual price
$\lambda^*_{\LAT}$ is positive. If instead the
binding resource is $\VRAM$, $\lambda^*_{\VRAM}$ is
positive. The planner's dual-price output
(Section \ref{sec:lagrangian}) answers the question
directly without profiling runs.

\emph{(b) What happens as $K \to \infty$?} Examine the
$\LAT$ scaling: $\alpha$ is $\Theta(K^2 M)$,
$\beta$ is $\Theta(KM + K^2 \lceil M/w \rceil)$,
$\gamma^*$ is $\Theta(K n R_r + K^2 (R_r + M))$. At
fixed $n, M$: $\beta$ and $\gamma^*$ are both
asymptotically $\Theta(K^2)$; $\alpha$ is worse by a
factor of $M$. The planner picks $\beta$ or
$\gamma^*$ based on whichever has the smaller
constant at the binding dimension.

\emph{(c) What's the minimum hardware that executes
this at budget $B$?} This is the inverse LP of
Section \ref{sec:lagrangian}: given $(K, n, M)$,
find the smallest $B$ such that at least one plan
is admissible. Plan $\beta$ at Instance 2 requires
$B_{\VRAM} \geq 4096$, $B_{\SMEM} \geq 0$, $B_{\REG}
\geq 3$, $B_{\LAT} \geq 3.3 \cdot 10^4$. If any of
these is tightened below threshold, the answer is
infeasibility; otherwise $\beta$ is realizable.

\subsection{The feedback loop}
\label{app-f:feedback}

The claims above are \emph{predictions} under
MODEL-4's symbolic parameters. In practice:

\begin{enumerate}[nosep,label=(\arabic*)]
  \item \emph{Specify} a workload $(K, n, M)$ and a
    budget $B$.
  \item \emph{Plan}: the planner returns a Pareto
    portfolio and dual prices.
  \item \emph{Execute} one or more plans from the
    portfolio on a real device; collect measured
    $(\VRAM, \SMEM, \REG, \LAT)$ values.
  \item \emph{Compare} measured vs.\ predicted cost.
    If the measured value deviates, update the
    hardware profile $\pi$: correct the per-atom
    baseline or the per-production combinator rule
    that produced the mis-prediction.
  \item \emph{Replan} with the updated $\pi$. The
    portfolio may change; phase boundaries shift.
    Re-measure.
  \item \emph{Iterate} until the profile stabilizes
    (successive replans produce the same plan) or
    until a target engineering budget is met.
\end{enumerate}

This loop is not a degenerate case. It is the
planner's operational mode. MODEL-4 with invented
parameters is the loop at iteration 0 with a
guessed $\pi$; item G1 of the reading guide is
iteration 1 (profile fit against a reference
device); item G2 is the binary that closes
the loop by making measurement possible at all.

The feedback structure is itself a $\kappa$-evolution
in [DAF]'s sense: each iteration refines the
discrimination regime $\Delta_\pi$ (the admissibility
regime indexed by the current profile). As $\pi$
converges, $\Delta_\pi$ stabilizes; the operationalist
axiom (Definition \ref{def:exec-operationalist})
becomes a statement about the limit regime, and the
portfolio output converges to the true Pareto
frontier for the measured hardware.

The shift in value proposition is concrete:

\begin{itemize}[nosep]
  \item \emph{Fast-execution framing:} the planner
    outputs a plan; the plan runs faster than a
    hand-coded alternative; profit. This framing
    requires implementation and benchmark wins to
    have any weight --- items G1, G2 \emph{and} a
    downstream demonstration.
  \item \emph{Feedback-loop framing:} the planner
    outputs a plan together with a cost prediction
    and dual prices; running the plan produces a
    measurement; the deviation between prediction
    and measurement is itself a \emph{discriminator}
    on the current hardware profile, in the
    operationalist sense. Engineering investment
    targets the residues with the largest downstream
    impact ($\lambda^*_r$ multiplied by the
    plan-family's sensitivity to $r$). The
    framework becomes an instrument for
    architecture-informed algorithm selection, not
    a one-shot compiler.
\end{itemize}

In the feedback-loop framing, the paper's $\gmark$
marks on \textsf{H}-class claims are not unresolved
defects --- they are \emph{loop invariants at
iteration 0}. Each iteration converts $\gmark$ to
$\cmark$ on the claims whose residues that iteration
witnessed; the paper's verification map is a
\emph{state} of the loop, not an end-state.

\subsection{Relationship to [DAF]'s self-hosting}
\label{app-f:self-hosting}

[DAF]'s Theorem 9.14 (Self-hosting) states that the
framework can serve as its own substrate: $\GF(2)
\to$ Belnap $\to$ Boolean dependent type $\to
\GF(2) \to \cdots$. The feedback loop described
above is a specific instance of this: the framework
specifies a $\kappa$ on plans, that $\kappa$ is
measured against hardware, the measurement updates
$\kappa$, and the loop continues.

Each iteration is an evolution in [DAF]'s sense; the
fixed point of the loop (if it exists) is a
hardware-calibrated planner whose $\gmark$ marks
have all been converted to $\cmark$, without any
definition or proof having changed.  The paper's
current state is the pre-fixed-point approximation.

This is the strongest form of the
framework-as-feedback-loop claim: the paper,
implementation, calibration, and measurement are
not four independent artifacts but four stages of
a single $\kappa$-evolving discrimination system
whose limit is the operationally-stable planner.
